\section*{Introdução}
\addcontentsline{toc}{section}{Introdução}
\label{sec:introducao}

% (==>) Este é o fenômeno; esta é a questão; há evidência científica suficiente para justificar investigá-la

% ECONOMIA + INSTITUIÇÕES --> DECISÃO DE USO DA TERRA --> ESTOQUE AMBIENTAL
% Foco do Capítulo 1 (H1): como retornos econômicos esperados e o ambiente institucional condicionam a conversão da vegetação nativa.
% (A retroalimentação — vegetação × clima --> produtividade — constitui o objeto de H2 no Capítulo 2; este artigo investiga a formação e transformação do estoque).

% HIPÓTESE DA TESE INVESTIGADA NESTE ARTIGO (H1 e desdobramentos):
% ---
% H1: RETORNOS, INSTITUIÇÕES E CONVERSÃO
% Retorno econômico esperado + ambiente institucional --> decisão de uso da terra --> taxa de conversão de vegetação nativa
%
% H1-a: PREÇOS E CAPACIDADE PRODUTIVA LOCAL
% Choques favoráveis de preços agrícolas geram maior pressão de conversão onde a capacidade produtiva da terra é mais elevada.
%
% H1-b: CONDICIONAMENTO INSTITUCIONAL
% Regras, monitoramento, fiscalização e sanções condicionam a sensibilidade da taxa de conversão ao retorno econômico esperado.
%
% Heterogeneidade: respostas distintas entre atividades (lavoura vs. pastagem) e entre biomas ( Amazônia, e Pantanal).

% EIXOS CONCEITUAIS QUE CONECTAM O ARTIGO À TESE:
% ---
% ENDOGENEIDADE AMBIENTAL
% O meio ambiente é tratado como componente endógeno do processo econômico: seu estoque é transformado pelas decisões produtivas sob incentivos e instituições.
%
% GOVERNANÇA TERRITORIAL
% Instituições, regulação e mecanismos de implementação operam sobre agentes e territórios heterogêneos, condicionando a intensidade da resposta aos incentivos de mercado.
%
% INTERTEMPORALIDADE
% Decisões presentes de conversão de uso da terra alteram estoques ambientais acumulados e condicionam oportunidades e riscos econômicos futuros.

% ------------------------------------
%%% (I) CONTEXTO
% ------------------------------------
% Mapeamento bibliográfico de H1:
% - Dilema conservação vs. rentabilidade e determinantes da decisão privada: Angelsen (2010), Lubowski et al. (2008), Jung et al. (2022).
% - Implementação regulatória e enforcement local: Viallon et al. (2019), Assunção et al. (2023), Pacheco & Meyer (2022).
% - Preços relativos, custos e contexto territorial: Harding et al. (2021), Chomitz & Thomas (2003).
% - Reorganização produtiva vs. expansão de fronteira: Bragança (2018).
% - Mato Grosso como contexto empírico e diversidade de trajetórias: VanWey et al. (2013), Simões et al. (2020).

%% Conflito: preservar ou converter (dilema econômico e determinantes privados/institucionais da decisão).
A expansão agropecuária costuma ser analisada à luz do dilema entre conservar a vegetação nativa, promovendo sustentabilidade ambiental, ou converter a terra para explorar usos economicamente mais rentáveis \cite{Angelsen2010}. Essa representação, contudo, simplifica uma decisão que não depende apenas da aptidão agronômica da terra, expectativa de valorização do lote ou da rentabilidade esperada da produção, mas também é influenciada por fatores institucionais, jurídicos e biofísicos \cite{LubowskiEtAl2008}. A conversão da vegetação é uma decisão privada realizada sob preços, custos de produção, restrições de crédito, expectativas, liquidez, direitos de propriedade, regras ambientais, \textit{enforcement} e probabilidade de sanções pelo uso irregular ou insustentável da terra \cite{JungEtAl2022}. Assim, territórios com expectativas de retornos agrícolas semelhantes podem apresentar distintos níveis de conversão quando diferem na configuração e na efetividade das instituições que regulam o uso da terra \cite{alstonTitlesConflictLand1999}.

%% Heterogeneidade institucional e implementação local (regras formais vs efetividade da fiscalização/enforcement).
Isto posto, uma mesma regra formal pode produzir respostas diferentes conforme a atividade econômica, o custo de oportunidade da terra e as condições locais de fiscalização \cite{AssuncaoEtAl2023}. Sob essa perspectiva, a conversão ou conservação da vegetação pode ser tratada como resultado endógeno da interação entre os retornos econômicos dos usos alternativos da terra e os incentivos, custos e riscos associados ao ambiente institucional \cite{Braganca2018}.

% Ponte para H1: retornos econômicos, incentivos institucionais e evolução do estoque de vegetação nativa.
Nessa toada, mudanças nos preços relativos, na rentabilidade agrícola, nos preços fundiários ou nas condições de acesso aos mercados podem aumentar o retorno esperado da conversão, enquanto incentivos positivos, fiscalização, sanções e restrições legais podem alterar a resposta dos agentes \cite{HardingEtAl2021}. Ato contínuo, compreender a evolução do estoque de vegetação requer analisar esse mecanismo econômico-institucional, pelo qual retornos econômicos, preços, custos e instituições alteram os incentivos privados à conversão ou à conservação da vegetação \cite{Braganca2018}.

% Justificativa para H1
Enfatize-se ainda que a resposta a esses incentivos pode envolver tanto a conversão de vegetação nativa quanto a reorganização da produção em áreas já convertidas. Os retornos da conversão e os custos institucionais enfrentados pelos agentes podem variar entre atividades, territórios e biomas \cite{ViallonEtAl2019}. Uma mesma alteração de preços pode, assim, produzir respostas distintas de uso da terra conforme as condições econômicas e territoriais sobre as quais os incentivos incidem \cite{ChomitzThomas2003}. É necessário, portanto, investigar quando e com que intensidade variações nos preços dos produtos agropecuários se traduzem em conversão da vegetação nativa e como essa relação é condicionada pela capacidade produtiva local da terra e pelo ambiente institucional.

%% Mato Grosso como caso empírico: grande escala agropecuária, três biomas e diferenciais biofísicos e econômicos.
Nesse contexto, Mato Grosso constitui um caso empírico especialmente informativo para investigar como os retornos econômicos dos usos agropecuários e o ambiente institucional se relacionam com as decisões de conversão da vegetação nativa. O estado combina produção agropecuária em larga escala com três biomas -- Amazônia, Cerrado e Pantanal --, cujas peculiaridades ambientais e econômicas resultam em diferenças relevantes de vegetação, clima, disponibilidade hídrica, sistemas produtivos e custo de oportunidade da terra dentro do mesmo território \cite{VanWeyEtAl2013}. Por esse norte, a literatura documenta tanto a expansão e a intensificação agrícola quanto a diversidade das trajetórias de uso da terra em Mato Grosso \cite{SimoesEtAl2020}.

% ------------------------------------
%%% (II) PROBLEMA, OBJETIVOS E JUSTIFICATIVA
% ------------------------------------

%% Problema central e formulação da questão de pesquisa (H1 em Mato Grosso).
Diante desse quadro, o problema central deste artigo consiste em compreender como incentivos econômicos e institucionais moldam a conversão ou manutenção da vegetação em ambientes ecologicamente e economicamente heterogêneos. O \textbf{problema} pode ser sintetizado no seguinte questionamento:

\medskip
\noindent\emph{Como retornos econômicos e instituições condicionam as decisões de conversão ou manutenção da vegetação nativa em Mato Grosso, e em que medida essa relação varia conforme a capacidade produtiva local, as atividades agropecuárias e os biomas?}
\medskip

%% Desdobramento do problema nas duas dimensões empíricas de H1: (i) preços × capacidade produtiva; (ii) condicionamento institucional; (iii) heterogeneidade entre atividades e biomas.
Essa pergunta desdobra-se em três dimensões empiricamente relacionadas: (i) se mudanças nos preços dos produtos agropecuários e nos preços da terra se traduzem em maior conversão onde a capacidade produtiva para esses usos é mais elevada; (ii) se as condições institucionais modificam a intensidade dessa resposta; e (iii) se há variações na resposta conforme as características das atividades e biomas.

%% Objetivo geral e objetivos específicos (i)–(iv): conversão bruta, resposta a preços × capacidade, instituições e heterogeneidade.
O \textbf{objetivo geral} deste artigo é investigar como e com que intensidade os determinantes observáveis do retorno econômico esperado e o ambiente institucional se relacionam com as decisões de conversão da vegetação nativa em Mato Grosso. Os \textbf{objetivos específicos} são: (i) construir medidas de conversão bruta da vegetação nativa para agricultura e pastagem; (ii) examinar a resposta das taxas de conversão a variações nos preços dos produtos agropecuários em interação com a capacidade produtiva local predeterminada; (iii) avaliar o condicionamento institucional dessa resposta, explicitando os limites de interpretação causal; e (iv) verificar a heterogeneidade dessas respostas entre atividades agropecuárias e biomas.

%% Relevância teórica e eixos da tese: meio ambiente endógeno, dimensão territorial e intertemporal.
A relevância deste trabalho decorre, em termos teóricos, de tratar o \textbf{meio ambiente} como componente \textbf{endógeno} do processo econômico, cujo estoque é transformado por decisões de uso da terra condicionadas por incentivos econômicos e instituições. Essa relação envolve uma \textbf{dimensão territorial}, na medida em que incentivos e instituições podem gerar respostas distintas conforme as condições locais (biomas), e uma \textbf{dimensão intertemporal}, pois a conversão presente altera estoques ambientais e oportunidades econômicas futuras.

%% Relevância aplicada e contribuição para políticas: decisões privadas, fiscalização pública, governança e limites institucionais.
Em termos empíricos e aplicados, compreender essas relações pode: (a) orientar decisões de firmas e produtores sobre conversão, conservação e manejo da terra, considerando os incentivos econômicos e as restrições institucionais associados aos usos alternativos; (b) oferecer subsídios à discussão de mecanismos de mercado que busquem compatibilizar produção agropecuária e conservação da vegetação nativa; e (c) subsidiar governos e organizações do terceiro setor responsáveis por fiscalização, conservação e alocação de recursos ambientais, ao examinar como incentivos econômicos, condições institucionais e características ambientais locais se relacionam com respostas distintas de conversão entre territórios.

% ------------------------------------
%%% (III) HIPÓTESES
% ------------------------------------

%% H1 — Governança territorial e heterogeneidade espacial: investiga como incentivos econômicos e instituições condicionam as decisões de conversão ou preservação da vegetação, considerando diferenças entre territórios e biomas.

%% Enunciação da hipótese H1 da tese e de seus desdobramentos empíricos (H1-a e H1-b).
Dito isso, o artigo examina a seguinte hipótese e seus dois desdobramentos:


% -------------------------------------------------------------------------
\medskip
% Hipótese geral do artigo (H1):
% CLAIM: H1-C001
% Maiores retornos econômicos esperados dos usos agropecuários elevam a taxa de conversão de vegetação nativa.

\noindent\textbf{(H1) Retornos, instituições e conversão.}

\noindent\emph{A taxa de conversão de vegetação nativa para usos agropecuários aumenta quando se eleva o retorno econômico esperado desses usos.}

% NOTA ECONOMÉTRICA — H1
%
% Variável dependente esperada:
%     - Conv_it = taxa de conversão de vegetação nativa para usos agropecuários no município/território i e período t.
%
% Construção candidata:
% Conv_it = (área de vegetação nativa em t-1 convertida para uso agropecuário em t) / (área de vegetação nativa disponível em t-1).
%
% Fonte candidata para Conv_it: MapBiomas — cobertura e transições anuais de uso da terra.
%
% Variável independente conceitual: R^e_i,t-1 = retorno econômico esperado dos usos agropecuários.
%
% R^e não precisa ser diretamente observado. Pode ser aproximado por determinantes observáveis do retorno esperado, especialmente:
%     - preços dos produtos agropecuários (World Bank Pink Sheet; CEPEA; IMEA; CONAB);
%     - capacidade produtiva/aptidão/produtividade potencial da terra (FAO-GAEZ; eventualmente IBGE/PAM em especificações alternativas);
%     - custos de produção/comercialização, caso haja dados adequados.
%
% Modelo-base esperado: Conv_it = beta_1 R^e_i,t-1 + alpha_i + lambda_t + gamma' Z_it + epsilon_it
%
% onde:
%     - alpha_i  = efeitos fixos territoriais/municipais;
%     - lambda_t = efeitos fixos de período;
%     - Z_it     = controles variantes no tempo;
%     - epsilon_it = termo de erro.
%
% Predição de H1: beta_1 > 0.
%
% Temporalidade inicial esperada: retorno em t-1 -> conversão em t.
% Testar posteriormente lags adicionais e efeitos acumulados.

%% Delimitação conceitual e variável dependente: vegetação nativa, conversão agropecuária e taxa de conversão.
Nesta hipótese, vegetação nativa corresponde à cobertura vegetal natural existente no território, enquanto sua conversão designa a substituição dessa cobertura por usos agropecuários, como lavouras ou pastagens. A taxa de conversão refere-se à parcela da vegetação nativa disponível em determinado território e período que é convertida para esses usos, constituindo o principal resultado de interesse. A ausência de conversão observada tampouco equivale, necessariamente, a uma decisão ativa de conservação ou à integridade ecológica da vegetação remanescente.

% -------------------------------------------------------------------------
\medskip
% Desdobramento empírico H1-a: interação entre choques de preços e capacidade produtiva predeterminada.
\noindent\textbf{(H1-a) Preços e capacidade produtiva local.}

\noindent\emph{O efeito de aumentos nos preços dos produtos agropecuários sobre a taxa de conversão de vegetação nativa é maior onde a capacidade produtiva local da terra para esses usos é mais elevada.}

% NOTA ECONOMÉTRICA — H1-a
%
% H1-a fornece uma operacionalização observável do retorno econômico esperado previsto em H1.
%
% Variável dependente:
% Conv_it = taxa de conversão definida em H1
%           (MapBiomas).
%
% Componente temporal:
% P_k,t = preço do produto agropecuário k no período t
%         (World Bank Pink Sheet; CEPEA; IMEA; CONAB).
%
% Transformações candidatas:
% ln(P_k,t), Delta ln(P_k,t), variação percentual, médias ou lags.
% A forma definitiva dependerá da estratégia de identificação.
%
% Componente espacial/predeterminado:
% S_ik = capacidade produtiva local da terra no território i
%        para o produto/atividade k.
%
% Proxies candidatas para S_ik:
% - aptidão agroecológica por cultura (FAO-GAEZ);
% - produtividade potencial por cultura (FAO-GAEZ);
% - produtividade histórica pré-determinada (IBGE/PAM);
% - composição produtiva em período-base (IBGE/PAM);
% - efetivo/estrutura pecuária prévia, se aplicável (IBGE/PPM).
%
% Construção candidata do choque de retorno:
%
% ShockReturn_it =
%     sum_k [ Delta ln(P_k,t-1) * S_ik ].
%
% Modelo esperado:
%
% Conv_it = beta_1 ShockReturn_it
%           + alpha_i + lambda_t
%           + gamma' Z_it + epsilon_it.
%
% Predição de H1-a:
% beta_1 > 0.
%
% Interpretação:
% um mesmo aumento de preço deve produzir maior pressão sobre a conversão
% onde a terra possui maior capacidade produtiva para a atividade
% beneficiada pelo choque.
%
% Possíveis desfechos adicionais:
%
% ConvAgric_it =
%   (vegetação nativa t-1 -> agricultura t)
%   / (vegetação nativa disponível t-1)
%
% ConvPast_it =
%   (vegetação nativa t-1 -> pastagem t)
%   / (vegetação nativa disponível t-1)
%
% Fonte:
% MapBiomas.
%
% Esses desfechos permitem separar respostas associadas a lavouras
% e à pecuária, sem transformar essa heterogeneidade em nova hipótese.

%% Identificação do retorno esperado: aproximação por determinantes observáveis (preços e capacidade produtiva como implicação observável de H1).
O retorno econômico esperado dos usos agropecuários corresponde ao benefício privado esperado da utilização produtiva da terra, considerando receitas e custos associados à atividade. Empiricamente, esse retorno não é diretamente observado nas bases reunidas, podendo ser aproximado por seus determinantes econômicos observáveis. Entre eles, destacam-se os preços dos produtos agropecuários e a capacidade produtiva local da terra. A H1-a expressa a expectativa de que choques favoráveis de preços produzam maior pressão sobre a conversão onde a terra oferece maior capacidade produtiva para os usos beneficiados. O teste dessa previsão será, portanto, um teste de uma implicação observável de H1, condicionado à validade da aproximação; não uma estimativa direta do retorno líquido esperado.

% -------------------------------------------------------------------------
\medskip
% Desdobramento empírico H1-b: condicionamento institucional sobre a sensibilidade da conversão aos retornos.
\noindent\textbf{(H1-b) Condicionamento institucional.}

\noindent\emph{A resposta da taxa de conversão ao retorno econômico esperado dos usos agropecuários é condicionada pelo ambiente institucional, particularmente pelas regras, pelo monitoramento, pela fiscalização e pelas sanções que alteram os custos e os riscos associados à conversão.}

%% Ambiente institucional: regras formais e mecanismos de implementação (monitoramento, fiscalização e sanções).
O ambiente institucional compreende as regras formais e os mecanismos de implementação que alteram os incentivos econômicos associados à conversão, incluindo direitos de propriedade, regulação ambiental, monitoramento, fiscalização e sanções. A H1-b não pressupõe que essas instituições determinem isoladamente a conversão, mas que modifiquem a intensidade com que alterações nos retornos econômicos esperados se traduzem em conversão efetiva. Quando uma medida representar, de maneira defensável, maior efetividade da fiscalização, espera-se menor sensibilidade da conversão ao retorno. Contagens de autos ou valores de multas, por si sós, não satisfazem essa condição.

%% Heterogeneidade empírica (atividades e biomas), escopo da análise e critérios para resultados nulos ou contrários.
As relações previstas em H1, H1-a e H1-b podem apresentar heterogeneidade entre atividades agropecuárias e biomas. Diferentes atividades apresentam tecnologias, estruturas de custos e retornos distintos, enquanto os biomas diferem em condições biofísicas e produtivas. Essa heterogeneidade será tratada como dimensão empírica da análise, sem constituir uma hipótese autônoma. Resultados nulos, de sinal contrário ou restritos a determinados grupos serão discutidos como limites ou evidência contrária às previsões propostas.

% ------------------------------------
%%% (IV) METODOLOGIA
% ------------------------------------

%% Estratégia empírica: dados em painel (2001–2024), fontes, ordenação temporal e identificação econométrica.
\textcolor{gray}{As hipóteses serão examinadas com dados em painel para Mato Grosso, tendo 2001--2024 como período candidato. A análise combinará transições anuais de uso da terra, preços agropecuários, medidas predeterminadas de capacidade produtiva e registros institucionais. As especificações considerarão a ordenação temporal entre incentivos e conversão, efeitos fixos territoriais e temporais e controles compatíveis com a identificação. A endogeneidade da fiscalização, a dependência espacial e a validade temporal das medidas de capacidade produtiva receberão tratamento explícito.}

% ------------------------------------
%%% (V) ORGANIZAÇÃO DO ARTIGO
% ------------------------------------

%% Organização do artigo: estrutura das seções subsequentes (1 a 5).
Além desta introdução, o artigo está organizado em cinco seções. A primeira revisa a literatura sobre retornos, instituições e conversão da vegetação. A segunda apresenta Mato Grosso como contexto empírico, descreve os dados e explicita a construção das variáveis. A terceira desenvolve a estratégia econométrica para examinar as hipóteses do artigo. A quarta estabelece os critérios de apresentação e discussão dos resultados a serem obtidos, com atenção à heterogeneidade e à robustez. Por fim, a quinta reúne as contribuições pretendidas e as limitações do estudo. Nesta versão para qualificação, não são apresentados resultados econométricos.

